\chapter{IPv6 Addressing}\label{ch:ipv6}
\bp{1.4}

\begin{outcomes}
By the end of this chapter you will be able to:
\begin{tight}
  \item \textbf{Abbreviate} and expand an IPv6 address correctly, applying both
        compression rules without ambiguity.
  \item \textbf{Identify} the address types the exam names: global unicast, link
        local, unique local, multicast, and the solicited-node address.
  \item \textbf{Construct} an interface identifier by modified EUI-64, which the
        blueprint calls out explicitly.
  \item \textbf{Size} a prefix, and \textbf{diagnose} the assignment faults that
        stop an IPv6 host communicating.
\end{tight}
\end{outcomes}

\begin{prereq}
\chref{ipv4} --- the ideas of network part, host part and prefix length carry
over unchanged. \chref{ethernet} for MAC addresses, which EUI-64 is built from.
\end{prereq}

\begin{keyterms}
global unicast $\bullet$ link local $\bullet$ unique local $\bullet$ multicast
$\bullet$ solicited-node $\bullet$ anycast $\bullet$ modified EUI-64 $\bullet$
SLAAC $\bullet$ router advertisement $\bullet$ neighbour discovery $\bullet$
dual stack $\bullet$ prefix
\end{keyterms}

\begin{examnote}[What v2.0 asks, and what it dropped]
Topic \tp{1.4} says \emph{troubleshoot} IPv6 address configuration, assignment
and prefix sizing, and narrows the scope to ``unicast and modified EUI~64''. The
previous blueprint asked you to compare every address type in detail; this one
does not. Concentrate on getting an address onto an interface, understanding
where the link-local one came from, and spotting a wrong prefix.
\end{examnote}

\section{Writing the address}

An IPv6 address is 128 bits, written as eight groups of four hex digits
separated by colons. Two rules compress it.

\begin{conceptbox}[The two compression rules]
\begin{tightnum}
  \item \textbf{Drop leading zeros} in any group. \cmd{0db8} becomes \cmd{db8};
        \cmd{0000} becomes \cmd{0}. Trailing zeros are never dropped.
  \item \textbf{Replace one run of all-zero groups with \cmd{::}}. You may do
        this \textbf{once only}, because a second \cmd{::} would make the address
        ambiguous --- there would be no way to know how many zeros each stood for.
\end{tightnum}
\end{conceptbox}

\begin{worked}[Compressing an address]
Start with the full form:

\cmd{2001:0db8:0000:0000:0000:ff00:0042:8329}

\begin{tightnum}
  \item Drop leading zeros in each group:
        \cmd{2001:db8:0:0:0:ff00:42:8329}
  \item Replace the longest run of zero groups with \cmd{::}:
        \cmd{2001:db8::ff00:42:8329}
\end{tightnum}

Expanding is the same in reverse. Given \cmd{2001:db8::1}, count the groups
present --- two before the \cmd{::} and one after, so three --- and the
\cmd{::} must therefore stand for five zero groups:
\cmd{2001:0db8:0000:0000:0000:0000:0000:0001}.
\end{worked}

\section{The types you must recognise}

\begin{longtable}{@{}L{3.0cm}L{3.4cm}L{8.8cm}@{}}
\toprule
\rowcolor{primary!15}
\textbf{Type} & \textbf{Prefix} & \textbf{What it is for} \\
\midrule
\endhead
Global unicast & \cmd{2000::/3} & Routable on the internet. In practice you meet
\cmd{2001:db8::/32}, which is reserved for documentation and used throughout
this book \\
\rowcolor{lightbg}
Link local & \cmd{fe80::/10} & Valid on one link only, never routed. Every IPv6
interface has one automatically, and routing protocols use it as the next hop \\
Unique local & \cmd{fc00::/7} & The rough equivalent of RFC 1918 private space.
In practice \cmd{fd00::/8} \\
\rowcolor{lightbg}
Multicast & \cmd{ff00::/8} & Groups. \cmd{ff02::1} is all nodes on the link,
\cmd{ff02::2} is all routers, \cmd{ff02::5} is OSPFv3 \\
Solicited-node & \cmd{ff02::1:ff00:0/104} & Used by neighbour discovery in place
of ARP. Formed from the last 24 bits of the unicast address \\
\rowcolor{lightbg}
Loopback & \cmd{::1} & The equivalent of \cmd{127.0.0.1} \\
Unspecified & \cmd{::} & ``No address yet'', used as a source during
autoconfiguration \\
\bottomrule
\end{longtable}

\begin{alertbox}[There is no broadcast in IPv6]
IPv6 has no broadcast address at all. Where IPv4 would broadcast, IPv6 uses a
multicast group --- usually \cmd{ff02::1}, all nodes. This is why ARP does not
exist in IPv6: its job is done by neighbour discovery using the solicited-node
multicast address, which reaches only the handful of hosts whose addresses end
the same way rather than every device on the segment.
\end{alertbox}

\section{Modified EUI-64}

The blueprint names this explicitly, so it is examinable arithmetic.

\begin{conceptbox}[Turning a 48-bit MAC into a 64-bit interface ID]
\begin{tightnum}
  \item Split the MAC in half.
  \item Insert \cmd{fffe} in the middle.
  \item \textbf{Flip the seventh bit} of the first byte --- the universal/local
        bit. In hex this means adding or subtracting 2 from the second hex digit.
\end{tightnum}
\end{conceptbox}

\begin{worked}[EUI-64 from 00d0.58a1.0201]
\begin{tightnum}
  \item Split: \cmd{00d0.58} and \cmd{a1.0201}.
  \item Insert \cmd{fffe}: \cmd{00d0:58ff:fea1:0201}.
  \item Flip the seventh bit of \cmd{00}. In binary \cmd{00} is
        \cmd{0000 0000}; flipping the seventh bit gives \cmd{0000 0010}, which
        is \cmd{02}.
  \item Interface ID is \cmd{02d0:58ff:fea1:0201}. With a prefix of
        \cmd{2001:db8:acad:10::/64} the full address is
        \cmd{2001:db8:acad:10:2d0:58ff:fea1:201}.
\end{tightnum}
The quick check: an address containing \cmd{ff:fe} in the middle of its host
portion was almost certainly built by EUI-64.
\end{worked}

\section{Configuring IPv6}

\begin{config}{IPv6 on a router interface, three ways}
! Routing must be enabled globally first -- this is the step most often
! forgotten, and without it the router will not forward or advertise.
ipv6 unicast-routing
!
interface GigabitEthernet0/0
 description User VLAN 10 gateway
 ! 1. A fully specified address.
 ipv6 address 2001:db8:acad:10::1/64
 ! 2. Or let the device build the host part by EUI-64.
 ipv6 address 2001:db8:acad:20::/64 eui-64
 ! 3. A link-local address is created automatically, but you can pin it
 !    to something memorable -- worth doing, because routing protocols
 !    show it as the next hop.
 ipv6 address fe80::1 link-local
 no shutdown
\end{config}

\begin{verify}{R1}{show ipv6 interface brief}
GigabitEthernet0/0     [up/up]
    FE80::1
    2001:DB8:ACAD:10::1
    2001:DB8:ACAD:20:2D0:58FF:FEA1:201
GigabitEthernet0/1     [administratively down/down]
    unassigned
\end{verify}

Note that one interface holds several addresses at once, and that this is normal
rather than a fault. An IPv6 interface essentially always has a link-local
address plus whatever global addresses it has been given.

\begin{conceptbox}[How a host gets an address: SLAAC versus DHCPv6]
A host sends a router solicitation to \cmd{ff02::2}; a router replies with a
\term{router advertisement} carrying the prefix. Then either:
\begin{tight}
  \item \textbf{SLAAC} --- the host builds its own address from the advertised
        prefix plus an interface ID it generates. No server involved.
  \item \textbf{DHCPv6} --- the advertisement tells the host to ask a server,
        which supplies the address and other options.
\end{tight}
A host with only a \cmd{fe80::} address and no global one has heard no router
advertisement. That is the IPv6 equivalent of an APIPA address, and it points at
the router, not the host.
\end{conceptbox}

\begin{pitfall}
\begin{tight}
  \item \textbf{Omitting \cmd{ipv6 unicast-routing}.} Interfaces come up and
        addresses look right, but the router neither forwards nor advertises.
  \item \textbf{Using \cmd{::} twice.} It is not merely discouraged, it is
        invalid, because the address can no longer be expanded uniquely.
  \item \textbf{Dropping trailing zeros.} \cmd{2001:db8:1000::} cannot be
        shortened to \cmd{2001:db8:1::}. Only \emph{leading} zeros go.
  \item \textbf{Assuming a prefix other than /64 for a LAN.} SLAAC requires a
        64-bit interface ID, so a LAN prefix longer than /64 breaks
        autoconfiguration. Point-to-point links are the usual exception.
\end{tight}
\end{pitfall}

\begin{breakfix}[Dual-stack host reaches IPv4 destinations only.]
A workstation has a working IPv4 address and, for IPv6, only
\cmd{fe80::2d0:58ff:fea1:201}. It cannot reach any IPv6 destination, on or off
the segment.

\textbf{Diagnosis.} A link-local address is generated by the host itself and
proves nothing except that the interface is up. The absence of a global address
means no router advertisement is being received. Either no router is configured
for IPv6 on that VLAN, or \cmd{ipv6 unicast-routing} is missing, or the router
interface has no global prefix to advertise.

\textbf{Fix.} On the gateway, confirm \cmd{ipv6 unicast-routing} is present in the
running configuration, then check \cmd{show ipv6 interface GigabitEthernet0/0}
shows a global address and that router advertisements are not suppressed. Bring
the prefix back and the host will build its own address within seconds --- no
change is needed on the host at all.
\end{breakfix}

\begin{lab}{Dual-stack a small network}{Packet Tracer}
\textbf{Topology.} One router, one switch, two PCs, addressed for both protocols.

\begin{tasks}
  \item Enable \cmd{ipv6 unicast-routing} and give \cmd{Gi0/0} the address
        \cmd{2001:db8:acad:10::1/64} and link-local \cmd{fe80::1}.
  \item Set both PCs to obtain IPv6 automatically. Record the addresses they
        build and confirm they contain \cmd{ff:fe}.
  \item Compute by hand what each PC's EUI-64 address should be from its MAC,
        and check your answer against the device.
  \item Ping between the PCs using IPv6, then ping the router's link-local
        address from a PC. Note that the link-local ping needs no global
        addressing at all.
  \item Remove \cmd{ipv6 unicast-routing} and observe what happens to the PCs
        after their addresses expire.
  \item \textbf{Extension.} Configure a second VLAN with prefix
        \cmd{2001:db8:acad:20::/64} and route between the two.
\end{tasks}
\end{lab}

\begin{checkpoint}
\begin{tightnum}
  \item Compress \cmd{2001:0db8:0000:0012:0000:0000:0000:00ab} fully.
  \item A host shows only an \cmd{fe80::} address. What is missing, and on which
        device would you look?
  \item Build the modified EUI-64 interface identifier for MAC
        \cmd{aabb.cc00.0100}.
\end{tightnum}
\end{checkpoint}

\begin{chaptersummary}
\begin{tight}
  \item Drop leading zeros; use \cmd{::} once, for the longest zero run.
  \item \cmd{2000::/3} global, \cmd{fe80::/10} link local, \cmd{fd00::/8} unique
        local, \cmd{ff00::/8} multicast. There is no broadcast.
  \item EUI-64: split the MAC, insert \cmd{fffe}, flip the seventh bit.
  \item \cmd{ipv6 unicast-routing} is a separate, easily forgotten step.
  \item LANs are /64 because SLAAC needs a 64-bit interface ID.
  \item Only a link-local address means no router advertisement was heard ---
        investigate the router, not the host.
\end{tight}
\end{chaptersummary}

\begin{reviewq}
  \item \textbf{(MCQ)} Which address is a valid compression of
        \cmd{2001:0db8:0000:0000:1000:0000:0000:0001}?
        (a)~2001:db8::1000::1 (b)~2001:db8:0:0:1::1 (c)~2001:db8::1000:0:0:1
        (d)~2001:db8:1000::1
  \item \textbf{(MCQ)} Which multicast group do all IPv6 routers on a link join?
        (a)~ff02::1 (b)~ff02::2 (c)~ff02::5 (d)~ff02::9
  \item \textbf{(Short)} Explain why \cmd{::} may appear only once in an address.
  \item \textbf{(Short)} A LAN is configured with a /112 prefix and hosts are set
        to autoconfigure. What will happen, and why?
  \item \textbf{(Applied)} You are dual-stacking an existing IPv4 network.
        Describe the order in which you would enable IPv6 on the router,
        verify each step, and what evidence you would collect before declaring
        the segment working.
\end{reviewq}
